% !TEX root = ../main.tex
\lecture{7}{2026-10-26}{Bewertung}

\section{Bewertung}

\begin{definition}
    [Bewertung] der Wert, der in der Bilanzpositionen steht.j
\end{definition}
\begin{itemize}
    \item Positionen
    \begin{itemize}
        \item Kasse
    \item Anlagevermögen
    \item 
\end{itemize}
\item Erstbewertung
\begin{itemize}
    \item Problemlos, weil Wert einer Maschine ihren Kaufpreis entspricht
\end{itemize}
\item Folgebewertung
\begin{itemize}
    \item Vermögenswerte werden abgeschrieben (devalued) \footnote{Was wäre der Prozess, der diese Abschreibung erklärt? Konstant Kapital usw nicht berücksichtigt }
    \item Abschreibung als Werkzeug zur Bewertung devaluierender Positionen 
\end{itemize}
\end{itemize}
\subsection{Grundproblematik}
Nächste Woche weil konzeptionell, heute nur technische Rechnung von 7.2 und .3
\subsection{Bewertung Aktiven}
\subsubsection{Methoden}
\begin{itemize}
    \item Lineare Abschreibungen (arithmetisch)
    \begin{itemize}
        \item Angenommen Anlage verliert jedes Jahr ein fixer Betrag an Wert bis 0
        \item Daraus wird durchschnittliche Abschreibung pro Jahr berechnet
        \item Abschreibungen in den spezifischen Perioden stimmen aber nicht zwingend 
    \end{itemize}
    \item Degressive Abschreibung (geometrisch)
    \begin{itemize}
        \item Verliert jedes Jahr Wert relativ zum aktuellen Wert
        \item Betragsmässig grössere Abschreibungen am Anfang
        \item Wird nicht zu null
    \end{itemize}
    \item Leistungsorientierte Abschreibung
    \begin{itemize}
        \item Anhand vom Output der Periode wird entsprechender Wert zu einem Verhältnis abgeschrieben 
    \end{itemize}
    \item Welche Methode soll man verwenden?
    \begin{itemize}
        \item Kommt auf die Anlage an
        \item Wenn es eine Anlage ist, die am Anfang des Gebrauchs bereits viel wert verliert, ist Degressive Methode zu bevorzugen
    \end{itemize}
\end{itemize}

\subsubsection{Darstellungsmöglichkeiten: Anlagevermögen}
\begin{itemize}
    \item Beispiel Fahrzeug
    \begin{itemize}
        \item Kaufpreis: 400k
        \item Abschreibung wird linear über 10 Jahre berechnet mit Rest 50k
        \item $\frac{400k - 50k}{10 Jahre} = 35k/Jahr$ 
    \end{itemize}
    \item Nettomethode (direkte Abschreibung)
    \begin{itemize}
        \item Bewertungen werden in einem entsprechendem Aufwand oder Ertrag ihre Spiegelung haben
        \item 1.1. Fahrzeuge / Bankguthaben 400k
        \item 31.12 Abschreibungen (Aufwand gestiegen) / Fahrzeuge (wird reduziert)  35k (die durchschnittlich berechnete Abschreibung pro Jahr)
    \end{itemize}
    \item Bruttomethode (indirekte Abschreibung)
    \begin{itemize}
        \item Mehr Arbeit notwendig für diese Methode, aber detaillierter
        \item 1.1. Fahrzeuge / Bankguthaben 400k
        \item 31.12. Abschreibungen / Wertberichtigungen Fahrzeuge (WB) 35k (Wert wird indirekt von Konto Fahrzeug abgeschrieben)
        \item Wertberichtigungen sind \textbf{negative Aktivkonten} auf aktiv Seite erfasst aber berechnet Abschreibungen
        \begin{itemize}
            \item Periodenüberschreitend kumuliert in WB-Konto 
        \end{itemize}
        \item In Bilanz
        \begin{itemize}
            \item Aktivseite: Fahrzeuge 400k
            \item - WB Fahrzeuge -35k (kumulierte Abschreibungen)   365k (zwischen Total)
        \end{itemize}
    \end{itemize}
\end{itemize}

Können Aktiven (zB Maschinen) ihre Werte erhöhen? Mehr Wert haben als am Anfang der Periode?
\begin{itemize}
    \item Werterhaltung
    \begin{itemize}
        \item z.B. Reparaturen
        \item sind Aufwände
        \item Reparaturaufwand / Kasse X (sofort in Erfolgsrechnung als Aufwand)
    \end{itemize}
    \item Wertvermehrung
    \begin{itemize}
        \item z.B. Modernisierungen
        \item Position in Bilanzkonto ("Aktivierung")
        \item Gebäude / Kasse X (nicht sofort in Erfolgsrechnung aber in kommenden Jahren )
    \end{itemize}
\end{itemize}

\subsubsection{Darstellungsmöglichkeiten: Umlaufvermögen}
\begin{itemize}
    \item Debitoren (Z.B wenn Rechnung nicht bezahlt wird)
    \item Darstellung?
    \begin{itemize}
        \item Konto '\textit{Delkredere}' Reduziert das Aktivkonto Debitoren \textbf{TODO: REVIEW}
        \item Keine Buchungen während des Jahres gemacht
    \end{itemize}
    \item Neues Erfolgsrechnungskonto: Debitorenverluste (Aufwand) mit 2 Komponente:
    \begin{itemize}
        \item effektive Zahlungsausfälle
        \begin{itemize}
            \item Debitorenverluste / Debitoren
        \end{itemize}
        \item Veränderungen des Kontos ''Delkredere''
        \begin{itemize}
            \item Erhöhung (Debitorenverluste /Delkredere)
            \item Reduktion (Delkredere / Debitorenverluste)
        \end{itemize}
    \end{itemize}
\end{itemize}
\subsection{Bewertung Passiven}
Grundsätzlich unproblematisch, weil im Beispiel einer Bankschuld 
Eigenkapital wird nicht Bewertet, weil es zu jedem Zeitpunkt aus der Differenz  der Aktiven und Verbindlichkeiten (FK) besteht.

\textbf{Ausnahme}: Rückstellungen  (z.B. Defekt verkaufter Ware)
\begin{itemize}
    \item Verbindlichkeiten, die mit Unsicherheit/Risiko behaftet sind
    \item ''nur'' hohe Wahrscheinlichkeit eines zukünftigen Mittelabflusses > 50\%
    \item z.B. Garantierückstellung oder Gerichtverurteilung
\end{itemize}

\textbf{Beispiel}
\begin{itemize}
    \item Garantierückstellung über Garantieaufwand
    \begin{itemize}
        \item Kaufpreis wird zurückgestattet
        \item 20
    \end{itemize}
\end{itemize}